\documentclass[11pt]{article}
\usepackage{geometry}
\usepackage{url}
\usepackage{color}
\geometry{left=1.3in,right=1.3in,top=1.5in,bottom=1in}

\begin{document}
\title{Response Memo}
\date{\today}
\maketitle

I would like to begin by thanking the editors for the invitation to revise and resubmit the manuscript to \emph{ASR}.  I would also like to sincerely thank the reviewers for their very incisive substantive criticism and very useful practical advice regarding the steps that I needed to take for an effective revision.  This were some of the most helpful reviews that I have ever received on a paper, which made the revision process a much smoother affair than it otherwise would have been.  

In what follows I detail the changes that I have made in this version of the paper in order to address the specific points of concern raised by all three readers.  I organize the discussion according to the three major areas that were highlighted by the editors in their letter.

\section*{Theoretical Claims}
\subsection*{Over-statement of the argument}
Reviewer A notes that the opening argument in the previous version of the paper was overstated. There I appear to imply that culture becomes ``socially meaningful'' \emph{only} insofar as it is used in interaction.  I am in complete agreement with this criticism.  It was certainly not my intent to make this (obviously unsupportable) suggestion.   I agree with the reviewer's proposal that most cultural consumption has a ``conspicuous'' aspect that is analytically and empirically distinct from the use of cultural knowledge in conversation.   In this version of the paper I have revised the wording and phrasing of the main argument so as to remove any remaining indication that mobilization in interaction is the only way in which culture plays an important social role. 

\subsection*{Better statement of research questions}
Reviewer A recommended that I add a transition paragraph between the end of the theoretical introduction and the beginning of the empirical analysis to indicate the research questions that will be addressed in the paper.  Reviewer C (point 1) notes a similar shortcoming with the previous version of the paper, in that the research questions are not entirely clear to the reader, a point that is highlighted by the editors.   

I found these observations to be very helpful.   I followed Reviewer A's editorial suggestion to include a section that clearly states that main research questions that will drive the empirical analysis up-front. The current version of the paper now contains a new section in which I outline and motivate the three main research questions that are the focus of the empirical analysis.  The inclusion of this section at the end of the theoretical introduction also serves to address a  concern noted by both Reviewer's A and C.  This regards the fact that the theoretically important discussion of intellectualization and the social construction of the symbolic value of cultural goods was given short-shrift in the previous version of the paper, since it was not brought up until the end.  In this version of the paper, this is introduced up-front, in order to motivate the third research question.

\subsection*{Conclusion}
Reviewer A noted with surprise that the concluding section appeared to be completely aimed at network researchers and should probably include a broader discussion of the implications for culture and stratification analysts.  Reviewer B also noted that this particular discussion appears to come out of the blue, and needs to be setup up-front.   I think these are two very good suggestions.  I have rewritten the conclusion to more clearly outline the broader contributions of the present research for culture and stratification analysts more broadly.  I have also attempted to anticipate the relevance of the findings for network analysis during the introduction. 

\section*{Measurement and Methods}
\subsection*{Addressing measurement issues early on}
Reviewer A: ``\emph{I would like to see a more careful assessment of the adequacy of the available measures with respect to the theory, and I would like to see at least some of this moved from the end of the paper (section 7.2) to the section in which the variables are introduced. In particular, the relative absence ofmeasures of `popular' culture consumption and talk while addressed in 7.2 needs to be confronted earlier, since otherwise it hangs over the whole paper.}''  This sentiment is echoed by Reviewer B, who also notes that the data limitations section should come earlier.  In this version of the paper I have moved the discussion of data limitations and measurement issues to the ``Data and Variables'' section as suggested by the reviewers.  I have also rewritten and expanded this section with an eye towards addressing the concerns expressed by reviewer A. 

\subsection*{Comparative validity of NSS survey}
Reviewer A suggested that I try to formally assess the relative validity of the NSS by comparing NSS distributions of some of the outcomes variables to those that can be obtained from other, better known surveys.  In this version of the paper I provide (footnote 4, page 16) some of these estimates (using the 2002 Survey for Public Participation in the Arts as the standard) for three such non-demographic variables:  frequency of attendance at galleries and museums, frequency of book reading and frequency of attendance at live performances.  The analysis suggest that NSS distributions of participants versus non-participants closely approximate those from the SPPA.  

\subsection*{Dichotomization of income variable}
Reviewer A asks why I dichotomized income and expressed surprised about the fact that I did not include the findings associated with income in the discussion. The short answer is that I initially was going to exclude income from the specification, but then decided to included in a truncated form.  I agree with the reviewer's point that if I include income as a right-hand predictor, then I should give it the same treatment as the other socio-demographic markers.  I have done that in this version of the paper.    I include the income full-range of income variables in all of the analysis reported in the paper; in addition, I discuss and compare the effects of income (in six categories) to the effects of the other socio-demographic factors.

\subsection*{More information on the data}
Reviewer A wondered about how many cases were dropped due to missing data. I include information on the original sample sizes (before list-wise deletion) on page 15.   Item non-response on the depedent variables never exceeded 1\% of the pooled sample.  The figures are shown in the table below.

\begin{table}[here]
\begin{tabular}{l c c}
\hline
\hline
&Consumption&Talk\\
\hline
Movies&0.33\% (6) &0.11\% (2)\\
Music&n/a&0.17\% (3) \\
Books&0.17\% (3) &0.17\% (3) \\
Arts&0.33\% (6) &0.11\% (2) \\
Performances&0.22\% (4) &0.06\% (1) \\
\hline
$N=1,803$
\end{tabular}
\end{table}

The reviewer also inquired about the response rate in the survey. Unfortunately, the technical documentation that accompanies the survey in the CPANDA archive only mentions the margin of error for the survey (3.1\%) but not the response rate.\footnote{\url{http://www.cpanda.org/cpanda/studies/a00229#method}.}

\subsection*{Statistical modelling issues}
Reviewer A had some concerns about statistical modeling.  In particular, the reviewer inquired as to whether there is any change in the parameter estimates when we compare singe-equation probit with bivariate probit models and as to whether these models are more sensitive to omitted variable bias.  In regards to the last issue, the models are no different from any other regression model; they are no more or less sensitive to specification issues than traditional methods.  In regards to the first issue, the answer is that the parameter estimates and standard errors change very little when go from the single-equation to the two equation model.  This can be seen in table~\ref{tab1}, where I compare the single-equation estimates for the effects of education, age, income and race on cultural talk involving books, with those obtained from the bivariate probit. 

\begin{table}[here]
\begin{tabular}{l*{2}{c}}
\hline
                    &         b/t&         b/t\\
				&   Single-Equation Probit & Bivariate Probit\\
educ==2         &0.1011 (0.9452)&0.1003 (0.9371)\\
educ==3         &0.4829 (4.4051)&0.4867 (4.4351)\\
educ==4         &0.5737 (4.9652)&0.5842 (5.0567)\\
educ==5         &0.6548 (5.1131)&0.6631 (5.1759)\\
age==2          &-0.1145 (-0.9563)&-0.1199 (-1.0039)\\
age==3          &-0.1856 (-1.5559)&-0.1896 (-1.5960)\\
age==4          &-0.2645 (-2.2331)&-0.2713 (-2.2938)\\
age==5          &-0.2207 (-2.1712)&-0.2286 (-2.2512)\\
age==6          &-0.2282 (-2.3417)&-0.2346 (-2.4102)\\
age==7          &-0.3757 (-3.6998)&-0.3786 (-3.7368)\\
gender==2       &0.4206 (7.8293)&0.4174 (7.7812)\\
finc==2         &0.0262 (0.2887)&0.0224 (0.2475)\\
finc==3         &0.0846 (0.9285)&0.0816 (0.8974)\\
finc==4         &0.2428 (2.4849)&0.2382 (2.4395)\\
finc==5         &0.0758 (0.6479)&0.0704 (0.6027)\\
finc==6         &0.2398 (1.9444)&0.2323 (1.8903)\\
race==2         &0.1604 (1.7683)&0.1634 (1.8035)\\
race==3         &-0.3979 (-1.8227)&-0.4040 (-1.8431)\\
race==4         &-0.0424 (-0.3302)&-0.0389 (-0.3034)\\
race==5         &0.3017 (1.4523)&0.2893 (1.3975)\\
chi2                &      180.58&      179.81\\
N                   &        1704&        1704\\
\hline
\end{tabular}
\label{tab1}
\end{table}

The advantages of the bivariate probit is that it allows us to get an estimate of the correlation between the joint outcomes, and that it allows us to test hypotheses about the relative equality of effects across equations in a straightforward way.  Beyond that, the bivariate probit provides essentially the same information about the relative effects of socio-demographic variables on each outcome as the single-equation version.   

The reviewer also wondered how much the parameter estimates change in the recursive model.   I present a relevant example in table~\ref{tab2}. Here the dependent variable is culture-talk involving music.  

\begin{table}[here]
\begin{tabular}{l*{2}{c}}
\hline
                &      b/t&      b/t\\
				&   Single-Equation Probit & Bivariate Probit\\
				&   & (Recursive) \\
educ==2         &0.0549 (0.3607)&-0.0603 (-0.4039)\\
educ==3         &0.4002 (2.5915)&0.0999 (0.6526)\\
educ==4         &0.4040 (2.4931)&0.0856 (0.5324)\\
educ==5         &0.5653 (3.0767)&0.1461 (0.7960)\\
age==2          &-0.1933 (-1.0344)&-0.2096 (-1.1401)\\
age==3          &-0.2726 (-1.3857)&-0.2571 (-1.3290)\\
age==4          &-0.4472 (-2.2659)&-0.3715 (-1.9102)\\
age==5          &-0.4321 (-2.6103)&-0.3670 (-2.2511)\\
age==6          &-0.4996 (-3.1849)&-0.5835 (-3.7933)\\
age==7          &-0.7064 (-4.3629)&-0.9277 (-5.8701)\\
gender==2       &0.0162 (0.1975)&-0.0511 (-0.6375)\\
finc==2         &0.1913 (1.3917)&0.2413 (1.8049)\\
finc==3         &0.1476 (1.0852)&0.2294 (1.7326)\\
finc==4         &0.2931 (2.0159)&0.3366 (2.3759)\\
finc==5         &0.0996 (0.5820)&0.0862 (0.5165)\\
finc==6         &0.3142 (1.7603)&0.3476 (1.9962)\\
race==2         &0.1564 (1.0993)&0.2615 (1.8625)\\
race==3         &0.1077 (0.3112)&0.1270 (0.3749)\\
race==4         &-0.3355 (-1.4789)&-0.3925 (-1.7958)\\
race==5         &0.3937 (1.5752)&0.2882 (1.1760)\\
favoperaclassical==1&0.2612 (2.4716)&1.6675 (13.7659)\\
chi2            &    68.74&   264.79\\
N               &      750&      750\\
\hline
\end{tabular}
\label{tab2}
\end{table}


In this case, there is a substantial difference between the estimates obtained from a single-equation specification and the bivariate probit specification.  First, the estimate for the coefficient for the recursive taste effect is driven \emph{downwards} ($\beta=0.26$ versus $\beta=1.67$)  in the single equation version (mainly because we are not accounting for the extraneous factors that affect both taste for classical music and rates of culture talk involving music, nor are we accounting for the fact that education, age and income also affect the likelihood of being a classical music aficionado.  This is what standard statistical theory leads us to expect).  

Second, in the single-equation version, the education effects appear to not change that much when we include taste for classical music.  In the bivariate probit model, however, we find that once we account for the effects of taste, then educational differences disappear ($|t|<1.96$).  This is consistent with the argument that taste for high-status culture is the \emph{intervening mechanism} that explains educational heterogeneity in culture-talk involving music.  We would have missed this if we did not account for the joint determination of taste and culture talk.  In this version of the paper, I attemtp to make this point in a more clear way, using the language of direct and indirect effects. 

In regards to the ``directionality'' issue, it is important to note that the recursive model does not settle the issue of causal direction in the strong sense.  These models are still subject to all of the limitations an caveats regarding causal inference that all cross-sectional, non-experimental research is subject to.  They are useful however, to get a clean estimate of the net effect of one variable on another variable in the presence of simultaneous determination of both variables by some other set of  unobserved factors and in light of the fact that both variables may themselves be determined by the same set of observed factors (as in the example above).   There is no (statistical) reason why a model in which a recursive effect going \emph{from} talk \emph{to} taste cannot be specified, but of course this model would make much less substantive sense and would violate the natural chronological order of causation.\footnote{The reviewer also asked that I elaborate and clarify some of these methodological issues in the paper.  In this version I have revised the introduction to the statistical models with a light of making some of these points clearer to the reader.}

\section*{Style, accessibility and length}
\subsection*{Manuscript Length}
Both the editors and Reviewers B and C noted that the previous version of the manuscript was too long, and that the paper could benefit from some rewriting aimed at succinctness, omitting redundancies, eliminating jargon and removing terminological ambiguity.  I thank Reviewers B and C for their editorial suggestions in this regard. This version of the manuscript features a much more streamlined introduction, as well as a much shorter results section.  I edited the introduction with an eye toward removing long extraneous quotes and towards simplifying word-choice and shortening the phrasing.  I removed all of the block-quotes and most of the long quotes.  In particular the suggestion to take the previous ``summary'' as a model to rewrite the introduction proved very helpful.   The paper is now about 10 double-spaced pages (2000+ words) shorter than the original, and the introduction is about four-doubled spaced pages shorter. 

\subsection*{Terminological choice}
Reviewer B made two very good suggestions in regards to terminology.  One was to stay away using cultural capital as a referent to what I'm measuring (since the concept has come to mean many things to different people and is thus confusing), and the other to stick with the terms ``culture-talk'' and ``culture-consumption'' for the two primary outcomes.  These points of concern are also raised by reviewer C, regarding the use of ``loose'' terminology, in particular as it pertains to the notion of cultural capital.   I agree with the reviewers that using cultural capital interchangeably with other terms may create confusion, especially given the polysemy of the notion of cultural capital.  In this version of the paper I have made sure to limit myself to the terms ``culture-talk'' and ``culture-consumption'' throughout the bulk of the discussion and presentation of the results. 

\subsection*{Better title}
Reviewer B writes: ``\emph{[y]ou need a better title, I suggest: `Cultural Consumption, Culture-Based Talk, and Stratification: Why `Cultural Matters'Matter.' This title more accurately represents what you do in the paper, and it creates an immediate tie to the important points you make in the conclusion.}''  I think the title change is a very good suggestion.  The current version of the paper now sports a title that is very close to that suggested by the reviewer. 

\subsection*{Removing the `disjunctions' section}
Reviewer A recommended that I consider removing the sections ``on disjunctions between interaction and consumption.'' I think this was very good advice, both for the problematic issues that the reviewer mentions, and also because this section simply reiterates findings that I discuss early on  I think the presentation of the results proceeds more smoothly without this section in the paper.  I thank the reviewer for this editorial suggestion. 

\subsection*{Minor editorial suggestions, spelling, etc.}
I have made the effort to catch most of the spelling/grammar issues identified by the readers, in particular Reviewer B (last section of the review). I have also thoroughly proof-read and spell-checked the paper.  

\end{document}